\section{Introduction}
\label{sec:introduction}

Deep learning has transformed histopathological image analysis, enabling
automated tissue classification, tumour grading, and biomarker
prediction~\cite{huang2024surface,kuntz2021gastrointestinal,chaunzwa2021deep,hirra2021breast,wetstein2022deep}.
Yet these models typically demand large annotated datasets to perform robustly,
a critical bottleneck in pathology where each sample must be labelled by an
expert. The problem is most acute for rare cancer subtypes and in
resource-limited settings, and it is compounded by the substantial variability of
histological images arising from tissue deformation, staining differences, and
scanner artefacts.

At the root of this challenge is the \textit{morphological heterogeneity} of
tissue. Unlike natural images, histological patches contain irregular structures
that vary dramatically across patients, preparation protocols, and scanning
conditions, so that the same tissue type can appear very different from slide to
slide. This forces models to allocate capacity to learning invariance to
nuisance geometric variation instead of diagnostic features, degrading both data
efficiency and cross-institution generalisation. Geometric normalisation---warping
irregular tissue into a standardised domain before feature extraction---has been
shown to help~\cite{huang2024surface}, but existing approaches require expensive
per-image optimisation, apply the same hand-crafted transformation to every
tissue type regardless of content, and decouple normalisation from classification
so the two cannot be jointly optimised. Stain-normalisation
methods~\cite{macenko2009method,shaban2019staingan} address colour, but not
geometric, variability and are therefore complementary.

We introduce \textbf{Deep Conformal Alignment (DCA)}, a framework that makes
geometric normalisation \textit{learnable} and integrates it directly into
end-to-end neural-network training. Building on spatial transformer
networks~\cite{jaderberg2015spatial} and diffeomorphic
registration~\cite{balakrishnan2019voxelmorph}, DCA parameterises deformations as
velocity fields integrated by scaling-and-squaring, guaranteeing
topology-preserving (diffeomorphic) transformations, and constrains them with a
differentiable \textit{conformal-energy} loss derived from the Cauchy--Riemann
equations so that learned warps preserve local angles---and thus local tissue
structure---while normalising global morphology. Because the transformation is
predicted in a single forward pass, DCA avoids the per-image optimisation of fixed
geometric methods and runs in real time. Methodological detail, including the
connection to quasiconformal mapping theory through the Beltrami coefficient, is
deferred to Methods so that the clinical results remain in the foreground.

Evaluating DCA on 100{,}000 colorectal-cancer patches across nine tissue classes,
we show that it (i) attains 93.2\% accuracy, outperforming spatial-transformer,
fixed-geometric, stain-normalisation, and self-supervised baselines with the same
backbone; (ii) is substantially more label-efficient, matching supervised
baselines that use several times more annotations; (iii) is complementary to
pathology foundation models, reaching 95.1\% accuracy when combined with a
pretrained encoder; (iv) generalises to an independent external cohort with the
smallest accuracy drop among compared methods; and (v) learns interpretable,
semantically adaptive transformations whose magnitude tracks tissue morphological
variability ($r=0.94$). Together these results indicate that domain-appropriate
geometric inductive bias is a practical route to data-efficient, robust
histopathology analysis.
